\documentclass[a4paper,10pt]{article}
\usepackage[utf8]{inputenc}
\usepackage[T1]{fontenc}
\usepackage[french]{babel}
\usepackage[top=1.6cm, bottom=1.6cm, left=1.5cm, right=1.5cm]{geometry}
\usepackage[table]{xcolor}
\usepackage{tabularx}
\usepackage{array}
\usepackage{enumitem}
\usepackage{titlesec}
\usepackage{hyperref}
\hypersetup{colorlinks=true, linkcolor=primary, urlcolor=primary}

\definecolor{primary}{HTML}{1A237E}
\definecolor{headbg}{HTML}{2C3E50}
\definecolor{cNav}{HTML}{1E88E5}   % Phase A navigation
\definecolor{cV2V}{HTML}{E65100}   % Phase B V2V
\definecolor{cP1}{HTML}{2E7D32}    % Phase C Protova1
\definecolor{cCons}{HTML}{6A1B9A}  % Phase D consolidation
\definecolor{rowalt}{HTML}{EEF2FA}

\titleformat{\section}{\large\bfseries\color{primary}}{}{0em}{}
\titlespacing*{\section}{0pt}{8pt}{3pt}
\setlength{\parindent}{0pt}
\renewcommand{\arraystretch}{1.3}

\newcommand{\g}[1]{\cellcolor{#1}}   % case coloree du Gantt

\title{\textbf{Plan de missions --- 2 mois}\\[2mm]
\large Projet PROTOVA (juillet--août 2026) --- ProtoVA1 \& ProtoVA2}
\author{Douae Sebti --- ALTEN Labs}
\date{3 juillet 2026}

\begin{document}
\maketitle
\vspace{-4pt}

\textbf{Objectif des 2 mois :} rendre \textbf{ProtoVA2 autonome de bout en bout}
(navigation Nav2), \textbf{enrichir et sécuriser le V2V}, et \textbf{remettre
ProtoVA1 en service} (odométrie encodeur $\rightarrow$ téléopération), pour
aboutir à une \textbf{démonstration intégrée} entre les deux véhicules.

% =====================================================================
\section*{\textcolor{primary}{Planning (diagramme de Gantt --- 8 semaines)}}
\vspace{2pt}
{\small
\begin{tabularx}{\textwidth}{|>{\raggedright\arraybackslash}X|c|c|c|c|c|c|c|c|}
\hline
\rowcolor{headbg}
\textcolor{white}{\textbf{Mission}} &
\textcolor{white}{S1} & \textcolor{white}{S2} & \textcolor{white}{S3} &
\textcolor{white}{S4} & \textcolor{white}{S5} & \textcolor{white}{S6} &
\textcolor{white}{S7} & \textcolor{white}{S8} \\
\hline
\multicolumn{9}{|l|}{\cellcolor{cNav!15}\textbf{\textcolor{cNav}{Phase A --- Navigation autonome ProtoVA2}}}\\
\hline
M1 : Cartographie propre (SLAM) + visualisation & \g{cNav} & & & & & & & \\
M2 : Convertisseur Ackermann (twist $\rightarrow$ servo) & \g{cNav} & \g{cNav} & & & & & & \\
M3 : Localisation AMCL + costmaps & & \g{cNav} & \g{cNav} & & & & & \\
M4 : Navigation Nav2 bout-en-bout + réglages & & & \g{cNav} & \g{cNav} & & & & \\
\hline
\multicolumn{9}{|l|}{\cellcolor{cV2V!15}\textbf{\textcolor{cV2V}{Phase B --- V2V enrichi et sécurisé}}}\\
\hline
M5 : Sécurité par certificats (Root/AA/ticket) & & & & \g{cV2V} & \g{cV2V} & & & \\
M6 : Messages DENM (freinage / obstacle) & & & & & \g{cV2V} & \g{cV2V} & & \\
\hline
\multicolumn{9}{|l|}{\cellcolor{cP1!15}\textbf{\textcolor{cP1}{Phase C --- ProtoVA1 remise en service}}}\\
\hline
M7 : Firmware Pico + node RX $\rightarrow$ odométrie & & & & & \g{cP1} & \g{cP1} & & \\
M8 : Intégration capteurs + téléopération & & & & & & \g{cP1} & \g{cP1} & \\
\hline
\multicolumn{9}{|l|}{\cellcolor{cCons!15}\textbf{\textcolor{cCons}{Phase D --- Intégration \& démonstration}}}\\
\hline
M9 : Démo V2V ProtoVA1 $\leftrightarrow$ ProtoVA2 + V2V$\rightarrow$AEB & & & & & & & \g{cCons} & \g{cCons} \\
M10 : Documentation + démonstration finale & & & & & & & & \g{cCons} \\
\hline
\end{tabularx}}

% =====================================================================
\section*{\textcolor{primary}{Détail des missions}}

\textbf{\textcolor{cNav}{Phase A --- Navigation autonome ProtoVA2 (S1--S4)}}
\begin{itemize}[leftmargin=1.4em, itemsep=1pt]
\item \textbf{M1 --- Cartographie propre.} Conduire pour cartographier le site
      (SLAM), sauver \texttt{.pgm + .yaml} ; corriger le transport WiFi des gros
      topics (profil FastDDS). \emph{Livrable :} carte exploitable + visualisation.
\item \textbf{M2 --- Convertisseur Ackermann.} N{\oe}ud convertissant la consigne
      Nav2 (twist) en angle de braquage servo pour le Pico
      ($\delta=\arctan(L\,\omega/v)$), testé roues en l'air. \emph{Livrable :}
      commande Nav2 $\rightarrow$ actionneurs opérationnelle.
\item \textbf{M3 --- Localisation AMCL + costmaps.} Configurer AMCL sur la carte,
      les costmaps (murs statiques + obstacles LiDAR live) et le contrôleur.
      \emph{Livrable :} le véhicule se localise et « voit » les obstacles.
\item \textbf{M4 --- Navigation bout-en-bout.} Envoyer un objectif depuis RViz,
      le véhicule s'y rend en évitant les obstacles ; réglage des vitesses et
      tolérances. \emph{Livrable :} \textbf{navigation autonome démontrée}.
\end{itemize}

\textbf{\textcolor{cV2V}{Phase B --- V2V enrichi et sécurisé (S4--S6)}}
\begin{itemize}[leftmargin=1.4em, itemsep=1pt]
\item \textbf{M5 --- Sécurité par certificats.} Générer la chaîne
      Root CA $\rightarrow$ AA $\rightarrow$ ticket (\texttt{certify}) et signer
      les CAM (\texttt{--security certs}). \emph{Livrable :} messages
      \textbf{signés et vérifiés} (vrai C-ITS).
\item \textbf{M6 --- Messages DENM.} Ajouter les alertes d'événement (freinage
      d'urgence, obstacle) émises à la demande. \emph{Livrable :} un véhicule
      freine $\rightarrow$ l'autre reçoit l'alerte DENM.
\end{itemize}

\textbf{\textcolor{cP1}{Phase C --- ProtoVA1 remise en service (S5--S7)}}
\begin{itemize}[leftmargin=1.4em, itemsep=1pt]
\item \textbf{M7 --- Odométrie (encodeur validé).} Monter l'encodeur, écrire le
      firmware Pico (Arduino : ticks + vitesse $\rightarrow$ ligne \texttt{MOTION})
      et le node RX (\texttt{/ticks}, \texttt{/velocity}). \emph{Livrable :}
      vitesse/odométrie réelles publiées sur ProtoVA1.
\item \textbf{M8 --- Capteurs + téléopération.} Intégrer LiDAR / caméra / IMU
      sous ROS 2 Foxy et rejouer la chaîne de téléopération (G29 / PS4).
      \emph{Livrable :} ProtoVA1 pilotable et instrumenté.
\end{itemize}

\textbf{\textcolor{cCons}{Phase D --- Intégration \& démonstration (S7--S8)}}
\begin{itemize}[leftmargin=1.4em, itemsep=1pt]
\item \textbf{M9 --- Démo V2V à deux véhicules.} Faire communiquer
      ProtoVA1 $\leftrightarrow$ ProtoVA2 (origine commune) et coupler l'alerte
      V2V au freinage d'urgence (AEB). \emph{Livrable :} scénario de sécurité
      coopérative réel.
\item \textbf{M10 --- Documentation + démonstration finale.} Rapport, présentation
      et démonstration intégrée. \emph{Livrable :} bilan de stage.
\end{itemize}

% =====================================================================
\section*{\textcolor{primary}{Jalons}}
\begin{itemize}[leftmargin=1.4em, itemsep=1pt]
\item \textbf{Fin S4 :} ProtoVA2 navigue de façon autonome (objectif RViz).
\item \textbf{Fin S6 :} V2V sécurisé (certificats) + DENM opérationnels.
\item \textbf{Fin S7 :} ProtoVA1 instrumenté et téléopérable.
\item \textbf{Fin S8 :} démonstration intégrée ProtoVA1 $\leftrightarrow$ ProtoVA2.
\end{itemize}

\vspace{4pt}
{\footnotesize\itshape Planning indicatif : les phases se recouvrent volontairement
(ex. V2V et ProtoVA1 en parallèle) ; les priorités pourront être ajustées selon les
aléas matériels (connecteur LiDAR, disponibilité des véhicules).}

\end{document}
